\documentclass[reqno]{amsart}
\usepackage[T1]{fontenc}
\usepackage{amsmath,amssymb,amsthm,mathtools}
\usepackage[shortlabels]{enumitem}
\usepackage{xr-hyper}
\usepackage[hidelinks]{hyperref}
\externaldocument[ft-]{20261124T213000--schwartz-space-cover-exercises__workbook}

\providecommand{\vol}{\operatorname{vol}}
\providecommand{\qr}[2]{\left(\frac{#1}{#2}\right)}
\newcommand{\SL}{\mathrm{SL}_2}
\newcommand{\tSL}{\widetilde{\mathrm{SL}}_2}
\newcommand{\Fc}{\mathcal{F}}
\newcommand{\Dc}{\mathcal{D}}
\newcommand{\Sc}{\mathcal{S}}
\newcommand{\Wc}{\mathcal{W}}
\newcommand{\Ab}{\mathbb{A}}
\newcommand{\oo}{\mathfrak{o}}
\newcommand{\pp}{\mathfrak{p}}

% Chapter number; \ref{chap} gives it.
\newcounter{chap}
\renewcommand{\thechap}{2}
\renewcommand{\thesection}{\thechap.\arabic{section}}
\numberwithin{equation}{section}
\theoremstyle{definition}
\newtheorem{problem}{Problem}[section]
\theoremstyle{plain}
\newtheorem{theorem}{Theorem}[section]

\newenvironment{hints}{\begin{description}[font=\normalfont\bfseries]}{\end{description}}
\newenvironment{exercises}{\subsection*{Exercises}\begin{enumerate}[1.]}{\end{enumerate}}

\setlength{\emergencystretch}{3em}
\title[The Schwartz space of the $n$-fold cover]{The Schwartz space of the
  $n$-fold cover of $\mathrm{SL}_2$: symmetries, the basic vector and
  inversion}
\author{}
\date{November 24, 2026}

\begin{document}
\maketitle
\refstepcounter{chap}\label{chap}

The previous sheet defined a Fourier transform $\Fc$ on genuine functions over the punctured plane, and the Schwartz space $\Sc := \Dc + \Fc \Dc$.  Is $\Fc$ an automorphism of $\Sc$, and does it fix the basic vector $\Phi^0$?  This sheet answers both questions, following \cite{N}.

\paragraph{Prerequisites.}
The sheet \emph{From Poisson summation on the plane to a Fourier transform for the $n$-fold cover of $\mathrm{SL}_2$}, whose notation and conventions we keep; it is Chapter~\ref{ft-chap}, and references such as \eqref{ft-eq:growth}, Theorem~\ref{ft-th:lines} or property~\ref{ft-h:tame} of the Hilbert symbol point there.  Fourier analysis on a local field.  Deligne's account of the cover \cite{Del} is the source of the formulas of \S\ref{sec:metaplectic}, but is not needed to do the problems.



%======================================================================
\section{The metaplectic group in the two charts}\label{sec:metaplectic}

Deligne lifts the action of $\SL(F)$ on $U$ to an action of Kubota's cover on the torsor \cite[\S\S1.7, 1.9, 1.11--1.12, 3.7]{Del}, and the lift of each unipotent matrix preserves one of the two trivializations \cite[(1.14.3)]{Del}.  In the charts, the lifts $n(u)$ and $n^-(u)$ of the upper and lower unipotent matrices with entry $u$ are
\begin{equation*}
  (n(u) \Phi)_X(x, y) := \Phi_X(x, y + u x),
  \qquad
  (n^-(u) \Phi)_Y(x, y) := \Phi_Y(x + u y, y),
\end{equation*}
the other component being determined by \eqref{ft-eq:transition} and continuity.  The group $\tSL(F)$ of operators generated by $\mu_n$ and the $n(u)$, $n^-(u)$ is a central extension of $\SL(F)$ by $\mu_n$; let $\pi$ be the projection \cite[\S\S1.11--1.12, 3.7]{Del}, \cite[Proposition~5.1]{N}.

\begin{problem}\label{pr:unip-other-chart}
  $(n(u) \Phi)_Y(x, y) = \{x, 1 + u x / y\}_n \Phi_Y(x, y + u x)$ when $x y (y + u x) \neq 0$, and $(n(u) \Phi)_Y(0, y) = \Phi_Y(0, y)$.
\end{problem}

\begin{problem}\label{pr:unip-lines}
  $R_X(n(u) \Phi)((x', y'), \tau) = R_X \Phi((x', y' + u x'), \tau)$.
\end{problem}

The analogues for $n^-(u)$ hold in the $Y$-chart.  Since a genuine function is determined by either component, Problem~\ref{pr:unip-lines} gives $\Fc n(u) \Phi = n(u) \Fc \Phi$ and $\Fc n^-(u) \Phi = n^-(u) \Fc \Phi$ on $\Dc$, and we obtain \cite[Corollary~5.3]{N}:

\begin{theorem}\label{th:metaplectic}
  Each $g \in \tSL(F)$ preserves $\Dc$ and $\Sc$, commutes with the dilations by $F^{\times n}$, satisfies $\Fc g \Phi = g \Fc \Phi$ for $\Phi \in \Dc$, and satisfies $\lvert g \Phi \rvert(v) = \lvert \Phi \rvert(v \, \pi(g))$.
\end{theorem}

Following Deligne \cite[\S1.9]{Del}, let $w(a) := n(a) n^-(-1/a) n(a)$ and $h(a) := w(a) w(1)^{-1}$, a lift of $\operatorname{diag}(a, a^{-1})$.  Deligne's lifts $d_1(a)$ and $d_2(a)$ of $\operatorname{diag}(a, 1)$ and $\operatorname{diag}(1, a)$ satisfy $(d_1(a) \Psi)_X(x, y) = \Psi_X(a x, y)$, $(d_2(a) \Psi)_Y(x, y) = \Psi_Y(x, a y)$, and $h(a) = d_1(a) d_2(a^{-1})$ \cite[\S1.12, (1.14.3), (1.9.3)]{Del}.

\begin{problem}\label{pr:torus-formula}
  For $a \in F^\times$,
  \begin{align*}
    (h(a) \Psi)_X(x, y) &= \{x, a\}_n \{-1, a\}_n \Psi_X(a x, y / a), \\
    (h(a) \Psi)_Y(x, y) &= \{a, y\}_n \Psi_Y(a x, y / a).
  \end{align*}
\end{problem}

The multiplier of $h(a)$ in the $X$-chart depends only on $x$, so $n(u) h(a) = h(a) n(u / a^2)$, and likewise $n^-(u) h(a) = h(a) n^-(a^2 u)$ \cite[Lemma~5.4]{N}.

\begin{problem}\label{pr:torus-center}
  $h(-1)^2$ is multiplication by $\{-1, -1\}_n$.  For $F = \mathbb{R}$ and $n = 2$, $w(1)^4 = -1$.
\end{problem}

The shears $n(u)$, $n^-(u)$ with $u \in \oo$ preserve the primitive vectors of $\oo^2$.  If they fix $\Phi^0$, then by Theorem~\ref{th:metaplectic} they fix $\Fc \Phi^0$ as well, and since they carry every vector to the line $x = 0$, a computation of $\Fc \Phi^0$ on that line determines it everywhere.

%======================================================================
\section{The basic vector}\label{sec:basic}

In this section $F$ is non-archimedean of residue characteristic prime to $n$.  We show that $\Phi^0$, defined in \S\ref{ft-sec:normalization} so that its Mellin coefficients behave correctly at one point, is fixed by $\Fc$.

\begin{problem}\label{pr:basic-invariant}
  $n(u)$ and $n^-(u)$ fix $\Phi^1$ and $\Phi^0$ for $u \in \oo$, and $h(a)$ fixes them for $a \in \oo^\times$.
\end{problem}

\begin{problem}\label{pr:axis-line}
  If $\Phi$ satisfies \eqref{ft-eq:growth} and $y', \tau \in F^\times$, then
  \begin{equation*}
    R_Y \Phi((0, y'), \tau) = \{-1, y'\}_n \lvert y' \rvert^{-1} \int_F \Phi_X(\tau^n / y', y) \, d y.
  \end{equation*}
\end{problem}

\begin{problem}\label{pr:basic-marginal}
  If $\psi$ is unramified, then $\int_F \Phi^0_X(x, y) \, d y$ is $1$ if $v(x) \in n \mathbb{Z}_{\geq 0}$, and $0$ otherwise.
\end{problem}

\begin{problem}\label{pr:basic-axis-values}
  If $\psi$ is unramified, then $(\Fc \Phi^0)_Y(0, y') = \Phi^0_Y(0, y')$ for $y' \in F^\times$.
\end{problem}

By Problem~\ref{pr:basic-invariant} and Theorem~\ref{th:metaplectic}, $\lvert \Fc \Phi^0 - \Phi^0 \rvert$ is invariant under the shears with $u \in \oo$, and these carry every point of $U$ to the line $x = 0$.  So Problem~\ref{pr:basic-axis-values} gives $\Fc \Phi^0 = \Phi^0$.

\begin{problem}\label{pr:basic-in-S}
  If $\psi$ is unramified, then $(1 - q^{-2}) \Phi^0 = \Fc \Phi^1 + q^{-n-1} \Phi^1(\varpi^n \, \cdot)$.
\end{problem}

Call $(F, \psi, n)$ \emph{good} if $F$ is non-archimedean of residue characteristic prime to $n$ and $\psi$ is trivial on $\oo$ \cite[Definition~6.4]{N}.  We have proved \cite[Proposition~6.3, Corollary~6.5]{N}:

\begin{theorem}\label{th:basic}
  The basic vector $\Phi^0$ lies in $\Sc$.  If $(F, \psi, n)$ is good and $b^{-1} \oo$ is the largest fractional ideal on which $\psi$ is trivial, then $\Fc \Phi^0 = \lvert b \rvert^n \Phi^0(b^n \, \cdot)$; in particular $\Fc \Phi^0 = \Phi^0$ if $\psi$ is unramified.
\end{theorem}

Problem~\ref{pr:basic-in-S} expresses $\Phi^0$ through $\Fc \Phi^1$, and a Poisson formula will need $\Fc$ to map $\Sc$ into itself.  For that we need $\Fc^2$ on $\Dc$.  For $n = 1$, $\Fc^2$ is computed by exchanging integrals along vertical lines with restriction to the line $x = 0$; we look for the same exchange.

%======================================================================
\section{Fourier inversion}\label{sec:inversion}

Let $F$ be a local field containing $\mu_n$.  For a genuine function $\Phi$ and $x, y \in F^\times$, let $m_\Phi(x) := \int_F \Phi_X(x, y) \, d y$ be its \emph{marginal} and $r_\Phi(y) := \Phi_Y(0, y)$ its \emph{axis values}.

\begin{problem}\label{pr:lin-marginals}
  For $n = 1$, $r_{\Fc \Phi}(y') = \int_F m_\Phi(x) \psi(x y') \, d x$ and $m_{\Fc \Phi}(x') = \int_F r_\Phi(y) \psi(- y x') \, d y$.
\end{problem}

For general $n$ the first relation persists.  Let $\Phi \in \Sc$, $y' \in F^\times$, and $\nu(\tau) := R_Y \Phi((0, y'), \tau)$.  By Problems~\ref{pr:axis-line} and~\ref{ft-pr:ft-root},
\begin{equation}\label{eq:marg-to-axis}
  \nu(\tau) = \{-1, y'\}_n \lvert y' \rvert^{-1} m_\Phi(\tau^n / y'),
  \qquad
  r_{\Fc \Phi}(\mu^n y') = \lvert n \rvert^{-1/2} \lvert \mu \rvert^{1-n} \hat{\nu}(\mu),
\end{equation}
for $\tau, \mu \in F^\times$ \cite[Lemma~A.1]{N}.  The second relation is the main local computation of \cite{N}; its proof rests on a twisted logarithm.  Let $f$ be a Schwartz--Bruhat function on $F$, $\eta$ a character of $F^\times$ of finite order, and
\begin{equation*}
  L(b) := \int_{F^\times} \eta(v) \bigl( f(b v) - f(v) \bigr) \, d^\times v
  \qquad (b \in F^\times).
\end{equation*}

\begin{problem}\label{pr:log-cocycle}
  $L(b c) = \eta(c)^{-1} L(b) + L(c)$.  If $\eta \neq 1$, then $L$ vanishes on the kernel of $\eta$.
\end{problem}

\begin{problem}\label{pr:steinberg-line}
  Let $s, y \in F^\times$ with $- s y \notin F^{\times n}$, and $\eta(v) := \{v, - s y\}_n$.  Then $L(\tau^n + s y) = 0$ for every $\tau \in F$ with $\tau^n + s y \neq 0$.
\end{problem}

\begin{problem}\label{pr:linear-factor}
  Let $\eta = 1$, $\rho \in C_c^\infty(F^\times)$ and $h := \hat{\rho}$.  For $a \in F$,
  \begin{equation*}
    \int_F h(\tau) L(\tau - a) \, d \tau = f(0) \int_{F^\times} \rho(\mu) \psi(\mu a) \, d^\times \mu.
  \end{equation*}
\end{problem}

When $- s y \in F^{\times n}$, the factorization $\tau^n + s y = \prod_{a^n = - s y} (\tau - a)$ and Problem~\ref{pr:linear-factor} apply; when $- s y \notin F^{\times n}$, Problem~\ref{pr:steinberg-line} applies.  With a bound $\int_{F^\times} \lvert f(b v) - f(v) \rvert \, d^\times v \ll 1 + \lvert \log \lvert b \rvert \rvert$ to justify the interchanges, this proves \cite[Lemma~A.2]{N}:

\begin{theorem}\label{th:axis-to-marg}
  Let $\Phi \in \Dc$ and $x' \in F^\times$, and let
  \begin{equation*}
    \sigma(\tau) := \lvert n \rvert \lvert x' \rvert^{-1} \lvert \tau \rvert^{n-1} r_\Phi(- \tau^n / x'),
    \qquad
    \sigma(0) := 0.
  \end{equation*}
  Then $m_{\Fc \Phi}(\mu^n x') = \lvert n \rvert^{-1/2} \hat{\sigma}(\mu)$ for $\mu \in F^\times$.
\end{theorem}

Let $\varepsilon := (-1)^{n+1}$.

\begin{problem}\label{pr:inversion-axis}
  For $\Phi \in \Dc$ and $y' \in F^\times$, $r_{\Fc \Fc \Phi}(y') = \{-1, y'\}_n r_\Phi(\varepsilon y')$, and this equals $r_{h(\varepsilon) \Phi}(y')$.
\end{problem}

Moving off the axis with $n^-(u)$, which commutes with $\Fc$ and with $h(\varepsilon)$, gives $\Fc^2 = h(\varepsilon)$ on $\Dc$ \cite[Proposition~A.3]{N}.  One sign in $\varepsilon = -(-1)^n$ comes from Fourier inversion in $\tau$, which replaces $\tau^n$ by $(-1)^n \tau^n$, and the other from the antisymmetry of $[\cdot, \cdot]$.  With the kernel identity $\overline{K(v, w)} = \{\varepsilon, x'\}_n K(\varepsilon w, v)$ one also obtains unitarity, and we arrive at \cite[Theorem~7.1]{N}:

\begin{theorem}\label{th:schwartz}
  The transform $\Fc$ maps $\Sc$ bijectively onto itself, with $\Fc^2 = h(\varepsilon)$; it commutes with $\tSL(F)$ on $\Sc$; and $\lVert \Fc \Phi \rVert_2 = \lVert \Phi \rVert_2$, where $\lVert \Phi \rVert_2^2 := \int_U \lvert \Phi \rvert^2$.
\end{theorem}

The element $h(\varepsilon)$ is central in $\tSL(F)$ and is the identity for odd $n$.  For $F = \mathbb{R}$ and $n = 2$ it exchanges opposite rays, so it need not act by a scalar on the functions satisfying \eqref{ft-eq:homogeneous}.

For $n = 1$, the Poisson formula contains $\Phi(0)$ and $\Fc \Phi(0)$.  For $n \geq 2$ and $\Phi \in \Dc$, Problem~\ref{ft-pr:schwartz-not-D} shows that $\Fc \Phi(\mu^n v)$ grows like $\lvert \mu \rvert^{1-n}$ times a function of $v$ as $\mu \to 0$.  There is no value at the origin, but there is a normalized limit along each orbit of the dilations; the sequel \emph{Leading terms and a Poisson formula for the $n$-fold cover of $\mathrm{SL}_2$} starts from it.

%======================================================================
\section*{Questions}

We do not know the answer to the following question.  Can $\Fc^2 = h(\varepsilon)$ be recovered by Mellin inversion from the Plancherel measures of Goldberg and Szpruch \cite{GS} and of Szpruch \cite{Sz}?


%======================================================================
\section*{Hints}

\subsection*{\S\ref{sec:metaplectic}}
\begin{hints}
\item[\ref{pr:unip-other-chart}] Apply \eqref{ft-eq:transition} at $(x, y)$ and at $(x, y + u x)$; near the axis, use property \ref{ft-h:open}.
\item[\ref{pr:unip-lines}] The shear, applied to $v$ and $v'$, preserves $[v, v']$, $x$ and $x'$.
\item[\ref{pr:torus-formula}] First compute the multiplier of $d_1(a)$ in the $Y$-chart from \eqref{ft-eq:transition}.
\item[\ref{pr:torus-center}] $w(-1) = w(1)^{-1}$, so $h(-1) = w(1)^{-2}$.
\end{hints}

\subsection*{\S\ref{sec:basic}}
\begin{hints}
\item[\ref{pr:basic-invariant}] For $\Phi^1$, the only issue is the factor $\{x, 1 + u x / y\}_n$ of Problem~\ref{pr:unip-other-chart}; use property \ref{ft-h:tame}.  For $\Phi^0$, the shears commute with dilations.
\item[\ref{pr:axis-line}] On this line $x = \tau^n / y'$, and $\{y, x\}_n = \{y, y'\}_n^{-1}$.
\item[\ref{pr:basic-marginal}] On the shells $v(y) = n j < v(x)$, the integrand is a power of the character $u \mapsto \{\varpi, u\}_n$.
\item[\ref{pr:basic-axis-values}] Problems~\ref{pr:axis-line} and~\ref{pr:basic-marginal} show that $R_Y \Phi^0((0, y'), \cdot)$ is a multiple of the indicator of a ball.
\item[\ref{pr:basic-in-S}] Apply $\Fc$ to Problem~\ref{ft-pr:basic-layer}, then dilate that identity.
\end{hints}

\subsection*{\S\ref{sec:inversion}}
\begin{hints}
\item[\ref{pr:lin-marginals}] The second relation needs Fourier inversion on $F$.
\item[\ref{pr:log-cocycle}] Substitute $v \mapsto v / c$, then exchange $b$ and $c$.
\item[\ref{pr:steinberg-line}] Divide both arguments of $\{\tau^n + s y, - s y\}_n$ by $\tau^n$ and use the Steinberg relation; then Problem~\ref{pr:log-cocycle}.
\item[\ref{pr:linear-factor}] Since $\int_F h = 0$, the term $f(v)$ drops out; then use Parseval and $\int_F \hat{f} = f(0)$.
\item[\ref{pr:inversion-axis}] Apply Theorem~\ref{th:axis-to-marg} with $x' = 1/y'$, then \eqref{eq:marg-to-axis} with $\mu = 1$, then Fourier inversion; for the last assertion use Problem~\ref{pr:torus-formula}.
\end{hints}

\begin{thebibliography}{McN}

\bibitem[Del]{Del}
P.~Deligne, \emph{Sommes de {G}auss cubiques et rev\^etements de {$\mathrm{SL}(2)$}, d'apr\`es {S}.~{J}. {P}atterson}, S\'eminaire Bourbaki, Exp.~539, Lecture Notes in Math., vol.~770, Springer, 1980, pp.~244--277.

\bibitem[GS]{GS}
D.~Goldberg and D.~Szpruch, \emph{Plancherel measures for coverings of {$p$}-adic {$\mathrm{SL}_2(F)$}}, Int. J. Number Theory \textbf{12} (2016), no.~7, 1907--1936.

\bibitem[N]{N}
P.~D. Nelson, \emph{A {S}chwartz space for the $n$-fold cover of $\mathrm{SL}_2$}, preprint, 2026.

\bibitem[Sz]{Sz}
D.~Szpruch, \emph{On {S}hahidi local coefficients matrix}, Manuscripta Math. \textbf{159} (2019), no.~1, 117--159.

\end{thebibliography}

\end{document}
